% Reconstructed from the compiled lecture-note PDF.
\chapter{Optional model: circular control and a conserved angular momentum}
\section{Why enlarge the triangle to a disk?}

The triangular control set has only discrete rotational symmetry. Replace it temporarily by a circular disk in the affine plane u0 + u1 + u2 = 1.


The modified problem is no longer exactly the Reinhardt packing problem, because some controls outside the triangle correspond to negative boundary curvature. It is a toy model. Its advantage is continuous rotational symmetry.


Continuous symmetries often produce conserved quantities. In elementary mechanics, invariance under spatial translation gives linear momentum, and invariance under rotation gives angular momentum. The same principle holds for optimal control.


\section{Noether’s principle in the present setting}

Let a one-parameter rotation act simultaneously on


\begin{displaytext}
X, Λ1, ΛR, Zu
\end{displaytext}

by conjugation. When the control set is a disk, the rotated control is still admissible, and the Hamiltonian is invariant. The control-theoretic Noether theorem then yields the conserved scalar


\begin{displaytext}
A = ⟨J, Λ1 + ΛR⟩.
\end{displaytext}

The source calls this angular momentum.


\begin{statementbox}{Theorem}
Theorem 13.1 (Angular momentum for circular control). Along every Pontryagin extremal of the circular-control problem, d ⟨J, Λ1 + ΛR⟩= 0. dt
\end{statementbox}

One can verify the conservation directly by differentiating and substituting the costate equations; the control-dependent terms cancel exactly because of rotational symmetry.


\section{The Cayley transform and hyperboloid coordinates}

\begin{displaytext}
The real algebra sl2(R) is isomorphic to \
( it z ! ) \
su(1, 1) = : t ∈R, z ∈C . \
¯z −it
\end{displaytext}

The isomorphism is a fixed change of basis called the Cayley transform. Its main advantage is that rotations act by multiplication of the complex coordinate by a phase.


For a determinant-one state, write


( q ) 1 + |w|2 w X = (−i q ). ¯w i 1 + |w|2


The pair q 1 + |w|2, w


lies on the upper sheet of the hyperboloid t2 −|w|2 = 1.


This is another standard model of the hyperbolic plane.


The other costates can be encoded by complex variables b, c. Near the singular locus,


\begin{displaytext}
w = b = c = 0.
\end{displaytext}

This is why these coordinates are adapted to the local analysis.


\section{Optimal control on a circle}

\begin{displaytext}
The boundary of the circular control set can be parametrized by a phase z with |z| = 1. In su(1, 1) \
coordinates, the control matrix has the form \
−iα βz ! \
Z∗(z) = , \
β/z iα
\end{displaytext}

after fixing positive constants α, β. The ratio


\begin{displaytext}
β \
ρ = \
α \
records the disk radius. The circumscribed disk of the triangle has ρ = 2; the inscribed disk has \
ρ = 1.
\end{displaytext}

\begin{displaytext}
Hamiltonian maximization and angular-momentum conservation reduce the optimal phase to a \
quadratic equation. In the source notation, if Q = [ΛR, X], then \
αQ21z2 −2iβQ11z + αQ12 = 0.
\end{displaytext}

Near the singular locus, the maximizing root satisfies c z∗= + O(|w|). |c|


Thus the optimal control direction is, to leading order, the phase of the third-chain variable c.


\section{Exact equations near the singular locus}

With fixed normal parameters


\begin{displaytext}
3 \
λcost = −1, d1 = 2, A = 3,
\end{displaytext}

the hyperboloid-coordinate equations imply


\begin{displaytext}
c \
w′ = −iρ + O(|w|), \
|c| \
b′ = 2iw + O(|b|), \
c′ = −3ib + O(|c| + |b| |w|2 + |b|2 |w|).
\end{displaytext}

If a trajectory reaches the singularity at t = 0, Gronwall estimates give


\begin{displaytext}
|w| = O(t), |b| = O(t2), |c| = O(t3).
\end{displaytext}

This justifies the Fuller truncation rigorously.


\section{Hamiltonian structure of the truncated system}

\begin{displaytext}
For \
z3 \
z′3 = z2, z′2 = z1, z′1 = −i |z3|,
\end{displaytext}

\begin{displaytext}
define \
i \
HF = 2(z2¯z1 −¯z2z1) + |z3| , \
AF = |z2|2 −(z1¯z3 + ¯z1z3).
\end{displaytext}

Both are conserved. The system is Hamiltonian for the nonstandard symplectic form


\begin{displaytext}
−1 1 ωF = dz1 ∧d¯z3 + dz2 ∧d¯z2 −1 dz3 ∧d¯z1. \
2i 2i 2i
\end{displaytext}

The unusual pairing of z1 with z3 reflects the third-order chain.


This Hamiltonian structure is conceptually valuable: the Fuller system is not an ad hoc approximation but retains the conservation laws and symplectic geometry of the original problem at leading order.


\section{Reducing by the virial symmetry}

\begin{displaytext}
On the zero set \
HF = AF = 0,
\end{displaytext}

introduce scale-invariant quantities


\begin{displaytext}
|z2| |z3| \
x2 = , x3 = . \
|z1|2 |z1|3
\end{displaytext}

The conservation laws imply


\begin{displaytext}
1 \
2 ≤x3 ≤x2. x2 > 0, x3 > 0, 2x2
\end{displaytext}

Thus (x2, x3) lies in a curved planar region Ω. Phase signs produce four copies of Ω, glued along their boundary curves to form a topological plane.


The Fuller flow descends, after a positive time rescaling, to an explicit vector field on this plane. It has two nontrivial equilibrium points:


\begin{displaytext}
2 r2 ! 2 r2 ! \
q∗+ = √ 10, 5, +, + , q∗−= √ 10, 5, −, − .
\end{displaytext}

\begin{displaytext}
They are the outward and inward logarithmic spirals. The Jacobian eigenvalues at q∗+ are \
√ √ \
− 2 ± i 3,
\end{displaytext}

so q∗+ is asymptotically stable; time reversal makes q∗−unstable.


\section{Global behavior of the circular Fuller model}

The source analyzes the vector field on the four glued copies of Ω. Apart from one special timereversal-symmetric trajectory passing through the corner (2, 2), the flow crosses the boundary curves in a controlled sequence. A Lyapunov function near q∗+ and monotonic winding around it show:


\begin{statementbox}{Theorem}
Theorem 13.2 (Global circular Fuller dynamics). Every trajectory other than the inward logarithmic spiral converges, after quotienting by scaling and rotation, to the outward logarithmic spiral as t →+∞. In reverse time, every trajectory other than the outward spiral converges to the inward spiral.
\end{statementbox}

This is a complete classification for the smooth circular model. It suggests the geometry later proved for the triangular Poincaré map: one inward fixed point, one outward fixed point, and a global basin for the outward one.


\section{What transfers to the triangular problem}

The following structural features survive when the circle is replaced by the triangle:


• weighted scaling of orders 1, 2, 3; • a time-reversal involution; • inward and outward self-similar spirals; • a stable outward angular profile and unstable inward profile; • global attraction on the blown-up angular space.


What changes is regularity. Circular control gives a smooth phase z∗. Triangular control jumps among three phases, so the continuous-time angular flow is replaced by a Poincaré map whose switching times are roots of cubic polynomials.


\begin{insightbox}{Chapter summary}
\end{insightbox}

\begin{displaytext}
Circularizing the control set creates continuous rotational symmetry and a conserved angular \
momentum. Hyperboloid coordinates center the singularity at w = b = c = 0 and rigorously \
produce a complex length-three Fuller system. After quotienting by scaling and rotation, \
that model has one stable outward spiral and one unstable inward spiral, with a global basin. \
This chapter supplies the conceptual template for the true triangular analysis, but the final \
proof does not depend logically on it.
\end{displaytext}
